% ════════════════════════════════════════════
% 模型的评价（7.模型评价.tex）写作规则 —— 模板内嵌，AI 先读本文件再写
% ════════════════════════════════════════════
\section{模型的评价}

\subsection{模型的优点}
\begin{itemize}[itemindent=2em]
\item \textbf{优点1：} 质量评价体系完整且可核验——逐指标给出 22 个质量指标的方向判定与语义依据（表\ref{tab:p1_dir}），其中争议最大的广告字段通过\textbf{原文抽查}确认为“第 1 类=无广告”的正向指标；对负向指标作补转换、对长度型指标转为中心接近度、对 8 个多维列表型指标用 softmax 期望压缩，逐字段审计出 19 处缺失并按中位数填充，并用抽样集与扩展集的域级质量对照（差值不超过 $0.002$）证明评分口径可跨集合迁移，从而避免了“方向未统一”这一常见错误。
\item \textbf{优点2：} 冲突分析给出了可比较的参照系——不仅报告冲突率（$22.3\%$）与 Kendall 协同系数（$0.069$），还计算了“指标独立”基准（$28.9\%$）并做阈值灵敏度检验（$\tau\in[0.5,2.5]$ 下比值恒 $<0.92$），从而客观得出“指标正相关但一致性不强”的结论，并归纳出三类冲突成因与消解规则，未回避数据中的不一致性。
\item \textbf{优点3：} 广义标度律的形式设计兼顾物理合理性与兼容性——采用质量缺口项 $C(1-Q)^{\gamma}$，在 $Q=1$ 时精确退化为经典 Chinchilla 形式，且不会在高质量端产生负损失（避免了 $E-CQ^{\gamma}$ 形式的非物理缺陷）；同时通过四种形式的同口径比较（$R^2$、AIC、内插与外推）择优，并用 B3 的独立插值轨迹验证数据项指数（$-0.2805$ vs $-0.2799$）。
\item \textbf{优点4：} 算力成本模型严格遵循赛题公式、不做简化——完整保留质量成本与训练 token 数成正比的结构（$C_Q=D[g(Q)-g(Q_0)]_+$）与附录 B 规定的三型函数参数，并解析给出临界上下文长度 $L_{\text{crit}}=6/\eta=30{,}000$ 与最优性条件式\eqref{eq:p3_balance}，用“等效附加数据量 $\tau(Q)$”把质量投入折算为 token 成本，使结构性转移具有清晰的解析解释。
\item \textbf{优点5：} 口径界定与证据分层规范——问题四明确给出时间轴、开源许可白名单与模型类型三套口径及筛选漏斗，完成附件 C8 的逐任务聚合（$1{,}860$ 个模型、$37$ 个子任务维度）并与 C1 做秩相关校验，用附件 C4 的算力增速与开源占比支撑情景假设；对 B2/B8 等口径异常的半合成数据主动作自洽性审计并划出可信度边界（B8 最小损失低于模型不可约损失下限）。
\end{itemize}

\subsection{模型的缺点}
\begin{itemize}[itemindent=2em]
\item \textbf{缺点1：} 配比—损失模型假设线性关系且未含领域交互项，13 个可测域的平均 $R^2$ 仅 $0.629$，说明仍有约 $37\%$ 的损失变异未被配比解释（可能来自随机种子、数据顺序等未观测因素）；nih\_exporter、enron\_emails、europarl、philpapers 四个无损失列的域只能固定于参考配比处理，其影响难以识别；追加混合质量指数后 $R^2$ 增量仅 $5\times10^{-5}$，虽有明确结论（质量信息已被配比吸收），但也意味着本文无法在线性框架内把 17 域质量评分进一步用于配比优化。
\item \textbf{缺点2：} 质量评价的指标方向仍部分依赖语义先验。广告、流畅度两个二分类字段已用原文与域对比完成实证核验，但 modernbert 系列的 6 级列表型指标、qurater 的 4 级评级等，“级别越高越好”仍来自字段语义而非独立验证；此外熵权法对分布形态敏感（左偏指标获更大权重），灵敏度检验显示改用等权平均后域级排序一致性为 $0.679$，方向判定若出错更会使排序发生实质变化（见第\ref{sec:check}节），说明域级质量排名的部分细节对赋权与方向设定敏感，引用时须同时给出其口径。
\item \textbf{缺点3：} 广义标度律的质量项依赖半合成数据标定（B6/B7），B8 因方向相反且损失低于不可约损失下限而不能使用，B2（Cerebras）的 MAPE 达 $33.7\%$，说明跨数据源的质量效应标定仍不充分；质量可行域受数据支持范围限制在 $[Q_0,1]$，向更低质量的外推缺乏数据支撑。
\item \textbf{缺点4：} 问题四的 Loss—Benchmark 桥接 $R^2$ 仅 $0.26$、MAPE 达 $67.8\%$，且 C1 中评测标准一致的窗口仅约 10 个月（2024-06 至 2025-03），2019--2023 的历史锚点口径不完全一致（已剔除 13 条异常条目），因此前沿预测的外推期约为一致窗口的 $1.2\sim2.4$ 倍，预测结果只能作为情景化区间；面板回归中规模与年份仍存在一定共线性，贡献占比对数据口径较敏感（长周期口径 $50.5\%$ vs 近端口径规模占比可正可负）。
\end{itemize}
